% !TeX spellcheck = en_US
\documentclass[aspectratio=1610]{beamer}
\mode<presentation>
{
  \usetheme{Boadilla}
  \setbeamercovered{transparent}
  
  \definecolor{cobalt}{rgb}{0.0, 0.20, 0.50}
  
  \usecolortheme[named=cobalt]{structure}
  %\usecolortheme{lily}
  %\setbeamercolor{headline}{fg=cyan!80!black}
  % \setbeamercolor{frametitle}{bg=white}.
}

\makeatletter
\define@key{beamerframe}{wide}[10pt]{%
	\def\beamer@cramped{\itemsep #1\topsep0.5pt\relax}}
\makeatother
	

\usepackage[english]{babel}
\usepackage[utf8]{inputenc}
\usepackage{times}
\usepackage[T1]{fontenc}
\usepackage{array}
\usepackage{xcolor}
\usepackage{colortbl}
\usepackage{booktabs}

\title[Fall - Lecture 10] % (optional, use only with long paper titles)
{Macroeconomic Theory and Policy: Lecture  10 (Ch.16)}

\author[University of Toronto] % (optional, use only with lots of authors)
{}
\date[] % (optional, should be abbreviation of conference name)

\begin{document}

\begin{frame}
  \titlepage
\end{frame}

%%%%%%%%%%%%%%%%%%%%%%%%%%%%%

\begin{frame}[wide]{What are You Going to Learn (Ch. 16)}
	\begin{itemize}
		\item The neoclassical consumption model
		\item Individuals choose the time path of their consumption to maximize utility
		
		\item Tow this model leads to a solution in which consumption is proportional to an individual’s total wealth
	\end{itemize}
\end{frame}

\begin{frame}[wide]{Introduction to Consumption Model}
	\begin{itemize}
		\item Consumption is the largest component of GDP (more than two-thirds)
		\item How do households decide what goods and services to consume?
		\vspace{1mm}
		\begin{itemize}
			\item In the Solow model, individuals spend a constant fraction of income  (i.e., $C_t = (1-\bar{s})Y_t$ )
			\item We should model how this $\bar{s}$ is formula, instead of taking it as exgenously given
		\end{itemize}
		\item Heterogeneity in consumer behavior at the microeconomic level
		\vspace{1mm}
		\begin{itemize}
			\item The rich follow the permanent-income hypothesis
			\item The poor typically have consumption that is quite sensitive to current income
		\end{itemize}
		\item The consumption and saving are interlinked
		\vspace{1mm}
		\begin{itemize}
			\item We observe a decline in the personal saving rate and the rise in the debt-income ratio in recent decades
		\end{itemize}
	\end{itemize}
\end{frame}

\begin{frame}[wide]{Introduction to Consumption Model}
	Neoclassical consumption model
	\vspace{2mm}
	\begin{itemize}
		\item Individuals choose consumption at each moment in time
		\item Objective: maximize a lifetime utility function 
		\vspace{1mm}
		\begin{itemize}
			\item Function depends on current and future consumption
		\end{itemize}
		\item People recognize that income in the future may differ from income today
		\vspace{1mm}
		\begin{itemize}
			\item Such differences influence consumption today
		\end{itemize}
	\end{itemize}
\end{frame}

\begin{frame}[wide]{The Neoclassical Consumption Model}
	\begin{itemize}
		\item Two time periods: today and the future
		\vspace{1mm}
		\begin{itemize}
			\item People earn income and consume in both periods
		\end{itemize}
		\item Key decision to make
		\vspace{1mm}
		\begin{itemize}
			\item How much to consume today versus how much to consume in the future 
		\end{itemize}
		\item The consumption model is based on two main elements:
		\vspace{1mm}
		\begin{itemize}
			\item An intertemporal budget constraint
			\begin{itemize}
				\item inter: between/across
				\item temporal: time
			\end{itemize}
			\item A utility function: a function that give us a measure of welfare
			
		\end{itemize}
	\end{itemize}
\end{frame}

\begin{frame}[wide]{The Intertemporal Budget Constraint}
	
	\begin{itemize}
		\item We start with two budget constraint, current period and future period
		\vspace{1mm}
		\begin{itemize}
			\item Each of them following that consumption equals income less saving
		\end{itemize}
		\item The consumer’s current-period budget constraint:
		\[c = y - (f'-f)\]
		where 
		\begin{itemize}
			\item[] $c$ is the consumption
			\item[] $f$ is the amount of wealth at the start of current period
			\item[] $f'$ is the amount of wealth at the start of future period
			\item[] $y$ is the amount of income
		\end{itemize}
		\item The $(f'-f)$ is the change in wealth = saving/borrowing
	\end{itemize}
\end{frame}

\begin{frame}[wide]{The Intertemporal Budget Constraint}
	\begin{itemize}
		\item In the last period (future) there is nothing to save for, so all wealth and income are used for consumption
		\item In general, we use the prime symbol " $'$ " to denote the variable in future period 
		\item The consumer’s future-period budget constraint:
		\[c' = y'  + (1+R)f',\]
		where 
		\begin{itemize}
			\item[] $R$ is the real interest rate
		\end{itemize}
	\end{itemize}
\end{frame}

\begin{frame}[wide]{The Intertemporal Budget Constraint}
	\begin{itemize}
		\item Arrange the future period budget constraint
		\[f' = \frac{c'-y'}{(1+R)}\]
		\item Then, sub it into the current period budget constraint
		\[c = y - \left(\frac{c'-y'}{1+R}-f\right)\]
		\item Put the consumption terms (both period) on the L.H.S. and the rest to the R.H.S
		\[c+\frac{c'}{(1+R)}  = y +f + \frac{y'}{(1+R)}\]
		%\item $\left(\frac{1}{1+r}\right)$ is the relative price of period 2 consumption goods in term of period 1 consumption (the price of period 1 consumption is one) 
	\end{itemize}
\end{frame}

\begin{frame}[wide]{The Intertemporal Budget Constraint}
	\begin{itemize}
		\item Human wealth is the present discounted value (PDV) of labor income
		\[\text{Present value of consumption} = \text{Financial wealth} + \text{Human wealth}\]
		\item Consumption in a given year can be different from income.
		\item PDV of consumption must equal lifetime resources
		\item Students anticipate an increase in future income and smooth their consumption
		\vspace{1mm}
		\begin{itemize}
			\item By borrowing against future labor income, college students can consume more today
			\item However, the present value of consumption must equal the present value of lifetime resources
		\end{itemize}
	\end{itemize}
\end{frame}

\begin{frame}[wide]{Utility}
	\begin{itemize}
		\item Utility depends on consumption today and in the future
		\item More consumption means higher utility 
		\[U = u(c) + \beta u(c')\]
		\item[] where $\beta$ is the discount factor or the weight placed on future consumption
		\item Diminishing Marginal Utility
		\begin{itemize}
			\item Each additional unit of consumption raises utility by a smaller amount
		\end{itemize}
		\item Usually, the $0<\beta<1$ so that future consumption is discounted -- Today’s consumption is valued more than future consumption
		
	\end{itemize}
\end{frame}

\begin{frame}[wide]{Flow Utility $u(c)$}
	\begin{itemize}
		\item A consumption level of c delivers a flow of utility to the consumer of $u(c)$
		\item Utility rises when $c$ increases, but the amount of the increase gets smaller and smaller, reflecting diminishing marginal utility
	\end{itemize}
	\begin{figure}
		\centering
		\includegraphics[width=0.4\linewidth]{"Figures/MACRO6_FIG16.01"}
	\end{figure}
	
\end{frame}

\begin{frame}[wide]{Choosing Consumption to Maximize Utility}
	\begin{itemize}
		\item Maximize utility by choosing current period consumption and future period consumption, subject to a budget constraint
		\[\max_{c,c'} U = u(c) + \beta u(c')\]
		s.t.
		\[c+\frac{c'}{(1+R)}  = y +f + \frac{y'}{(1+R)} = \bar{X} \]
		\item[] where $\bar{X}$ represents the total wealth
	\end{itemize}
	
\end{frame}

\begin{frame}[wide]{Choosing Consumption to Maximize Utility}
	\begin{itemize}
		\item Rearrange the budget constraint:
		\[c  =\bar{X} -\frac{c'}{(1+R)}\]
		\item Sub it into the utility function, the maximization problem becomes
		\[\max_{c'} U = u\left(\bar{X} -\frac{c'}{(1+R)}\right) + \beta u(c')\]
		\item First order conditions (take the derivative w.r.t. $c'$ and set it equal to zero):
		\[\frac{d u\left(\bar{X} -\frac{c'}{(1+R)}\right)}{d c} \frac{-1}{1+R} + \beta \frac{du(c')}{dc'}=0\]
		The first term here we use the chain rule of calculus. Let $u'(c') = \frac{du(c')}{dc'}$, etc
		\[\beta u'(c')= \frac{u'(c)}{1+R}\]
		
	\end{itemize}
\end{frame}

\begin{frame}[wide]{Choosing Consumption to Maximize Utility}
	\begin{itemize}
		\item We want to make ourselves as happy as possible, but have a limited income
		\item Marginal utility of today’s consumption: $u'(c)$
		\item Marginal utility of future consumption: $u'(c')$
		\item As seen from the equation from previous slides, agents can either 
		\begin{itemize}
			\item consume 1 unit today or 
			\item save that unit and consume $1+R$ units in the future
		\end{itemize}
		\item If individuals maximize utility, agents must be indifferent between consuming today or in the future
	\end{itemize}
\end{frame}

\begin{frame}[wide]{Choosing Consumption to Maximize Utility}
	\begin{itemize}
		\item The equation that govern the trade-off between current and future consumption is called Euler equation:
		\[\beta (1+R) u'(c')= u'(c)\]
		\item 	MU of consuming 1 more unit today = discounted consumption of  units in the future
		\begin{itemize}
			\item LHS: Saving that same unit and consuming 1+R units in the future gives extra utility (discounted)
			\item RHS: Consuming one unit today, the agent gets the MU of c today
			
		\end{itemize}
		\item Must be indifferent between consuming today and tomorrow
		\item If the MU of $c$ is greater than the utility from consuming the extra unit (plus interest) in the future,
		\begin{itemize}
			\item the individual should reduce consumption in the future and raise it today, increasing overall lifetime utility
			\item This reallocation should occur until the agent is just indifferent between consuming an extra unit today and consuming that unit plus interest in the future
		\end{itemize}
	\end{itemize}
\end{frame}

\begin{frame}[wide]{Solving the Euler Equation}
	\begin{itemize}
		\item Take $u(c) = log (c)$ as an example
		\begin{itemize}
			\item Recall that $log (c) = \frac{1}{c}$
		\end{itemize}
		\item This implies that $\beta (1+R) u'(c')= u'(c)$ would become
		
		\[\frac{\beta (1+R)}{c'} = \frac{1}{c}\]
		\item Rearranging the Euler equation:
		\[\frac{c'}{c} = \beta (1+R)\]
		\item Left side is growth rate of consumption (plus 1)
		\begin{itemize}
			\item Consumption is chosen such that the growth rate of consumption ($c'/c$) is a product of the discount parameter and the interest rate
		\end{itemize}
		\item Agent takes R as given in partial equilibrium consumption problem, choosing the consumption growth rate
	\end{itemize}
\end{frame}

\begin{frame}[wide]{Solving the Euler Equation}
	\[\frac{c'}{c} = \beta (1+R)\]
	\begin{itemize}
		\item Lower $\beta$: Impatient, consumption growth is lower
		\item Higher $R$: Patient, consumption growth is faster
		\item Suppose the economy is in steady state (i.e., $c'/c = 1$), then
		\[\beta= \frac{1}{1+R}\]
	\end{itemize}
\end{frame}

\begin{frame}[wide]{Solving the Model}
	\begin{itemize}
		\item Rearrange the Euler equation:
		\[c' = \beta (1+R)c\]
		\item Plug it into the budget constraint
		\[c+\frac{\beta (1+R)c}{(1+R)}  = y +f + \frac{y'}{(1+R)}\]
		\[\implies c^*  = \left[y +f + \frac{y'}{(1+R)}\right] \frac{1}{1+\beta} = \bar{X}\frac{1}{1+\beta}\]
		\item Plus $c^*$ into the Euler equation (alternative, you can plug it into the budget constraint)
		\[c^{'*} = \left[y +f + \frac{y'}{(1+R)}\right] \frac{\beta (1+R)}{1+\beta}= (1+R)\bar{X}\frac{\beta }{1+\beta}\]
	\end{itemize}
\end{frame}

\begin{frame}[wide]{The Effect of a Rise in $R$ on Consumption}
	\begin{itemize} 
		\item Wealth Effect: A higher interest rate will reduce this present value, therefore reducing consumption in the case of log utility
		\begin{itemize}
			\item Total wealth depends on the interest rate because of PDV of labor income (i.e., $\frac{y'}{1+R}$)
		\end{itemize}
		\item Substitution effect: When interest rate increase, current consumption become more expensive (because saving will lead to even more future consumption), so consumers will reduce current consumption today
		\item Income effect: consumers are now richer—because their current saving leads to more income in the future—which makes them want to consume more today
		\item With log utility, the IE and SE offset each other
		\item Consumption could increase or decrease, depending on utility function
	\end{itemize}
\end{frame}

\begin{frame}[wide]{Lessons from The Neoclassical Model}
	\begin{itemize}
		\item The permanent-income hypothesis is an implication of this model
		\begin{itemize}
			\item Consumption depends on average income (PDV of income) rather than current income 
		\end{itemize}
		\item Recall overall wealth: $\bar{X} = f + y + \frac{y'}{1+R}$
		\item[] Last two terms are PDV of income
		\item Consumption is proportional to a consumer’s overall wealth
	\end{itemize}
\end{frame}

\begin{frame}[wide]{Lessons from The Neoclassical Model}
	\begin{itemize}
		\item Permanent-income hypothesis implies that we want to smooth our consumption over time
		\item What happens if income increases?
		\begin{itemize}
			\item We may spend some of the extra income, but will save some for the future
		\end{itemize}
		\item Consumer do not like variability in consumption
		\item Marginal propensity to consume: portion of additional income that is spent for consumption
		\begin{itemize}
			\item In the log example, it is equal to $\frac{1}{1+\beta}$ for $c$ and $\frac{\beta}{1+\beta}$ for $c'$
		\end{itemize}
		
	\end{itemize}
\end{frame}

\begin{frame}[wide]{The Desire to Smooth Consumption}
	\begin{itemize}
		\item Suppose an individual could either consume $c_1$ today and $c_2$ in the future, or could consume the average of these two values in both periods
		\item Because of diminishing marginal utility, agents prefer to smooth consumption and take the average in both periods
		\item That is, lifetime utility is higher when consumer gets $u(\bar{c})$ in both periods, rather than $u(c1) + u(c2)$
	\end{itemize}
	\begin{figure}
		\centering
		\includegraphics[width=0.35\linewidth]{"Figures/MACRO6_FIG16.02"}
	\end{figure}
	
\end{frame}

\begin{frame}[wide]{Ricardian Equivalence}
	\begin{itemize}
		\item Consumption depends on PDV of taxes and is invariant to the timing of taxes
		\item The government has a similar constraint that the lifetime government spending is also equal to lifetime tax revenue
		\item A tax cut today financed by higher future taxes will not affect consumption, as long as the lifetime government spending is the same
		\[c+\frac{c'}{(1+R)}  = y -t +f + \frac{y'-t'}{(1+R)} = \bar{X}\]
		\begin{itemize}
			\item Think of there is an increase in $t$ with a decrease in $t'$ of the same size (apply the interest rate)
			\[c+\frac{c'}{(1+R)}  = y -(t+\Delta t) +f + \frac{y'-(t'-(1+R)\Delta t)}{(1+R)}\]
			\item A change in the timing of taxes does not change wealth, so consumption doesn’t change
		\end{itemize}
	\end{itemize}
\end{frame}

\begin{frame}[wide]{Borrowing Constraints}
	\begin{itemize}
		\item The model assumes that anyone can borrow or save at the interest rate $R$
		\item However, in reality, some people may not be able to borrow because of 
		\begin{itemize}
			\item poor financial wealth or no access to credit
			\item weak economic conditions as a whole 
		\end{itemize}
		\item In either of these cases, the intertemporal budget constraint is no longer the correct constraint
		\item The constraint on consumption for individuals with no financial wealth and no access to credit can be rewritten as
		\[c\leq y\]
		\item The inequality assumes people are constrained by the lack of borrowing opportunities to be no greater than a person’s income
	\end{itemize}
\end{frame}

\begin{frame}[wide]{Borrowing Constraints}
	\begin{itemize}
		\item If a consumer is already consuming less than their income, the no-borrowing constraint will be nonbinding
		\begin{itemize}
			\item The individual is already saving, so not allowing him or her to borrow does not change anything
		\end{itemize}
		\item If income is low and the borrowing constraint binds, then $c=y$
		\begin{itemize}
			\item In most binding cases, it is undesired outcome
			\item People cannot smooth their consumption even if they want to
		\end{itemize}
	\end{itemize}
\end{frame}

\begin{frame}[wide]{Case Study: Consumption versus Expenditure}
	\begin{itemize}
		\item Key prediction of model:	Consumption should not change when an anticipated event occurs
		\begin{itemize}
			\item Consumer wants to smooth out their consumption over time
		\end{itemize}
		\item However, this doesn’t always hold true
		\begin{itemize}
			\item A retirement is usually not an unexpected event
			\item Consumption should remain relatively unchanged when people retire
			\item However, expenditures on consumption change quite markedly around this event, falling by around 17 percent
		\end{itemize}
		\item A study by Aguiar and Hurst look at it differently
		\begin{itemize}
			\item Use data set of food diaries for a large number of households
			\item While expenditures on consumption do indeed decline sharply upon retirement, the food and calories that people actually consume show no such decline
			\item Instead, households spend much more time shopping for food and preparing it themselves
		\end{itemize}
	\end{itemize}
\end{frame}

\begin{frame}[wide]{Precautionary Saving}
	\begin{itemize}
		\item Saving money to hedge against the possibility of a drop in income due to uncertain and unpredictable event
		\begin{itemize}
			\item Consumers might save when income and wealth are temporarily low, meanwhile the basic permanent-income hypothesis would suggest borrowing
		\end{itemize}
		\item The recent financial crisis provides an excellent example of precautionary saving
		\begin{itemize}
			\item Saving rates rose sharply during the financial crisis, and precautionary motives are a logical part of the explanation
		\end{itemize}
	\end{itemize}
\end{frame}




%\begin{frame}[wide]{An Increase in the Investment Rate}
%	\begin{columns} 
%			% Column 1
%			\begin{column}{.5\textwidth}
%				\begin{itemize}
%					\item Suppose the investment rate increases permanently for exogenous reasons ($s'>s$)
%					\begin{itemize}
%						\item The investment curve rotates upward ($\bar{s}'\bar{A}  \bar{L}^{1-\alpha} K_t^\alpha > \bar{s}\bar{A}  \bar{L}^{1-\alpha} K_t^\alpha$)
%						\item The depreciation curve remains unchanged ($\delta K_t$)
%						\item The capital stock increases by transition dynamics to reach the new steady state
%						\begin{itemize}
%							\item This happens because investment exceeds depreciation ($sY_t>\delta K_t$)
%						\end{itemize}
%						\item The new steady state is located to the right
%					\end{itemize}
%				\end{itemize}
%			\end{column}
%			% Column 2    
%			\begin{column}{.5\textwidth}			
%				\begin{figure}
%					\centering
%					\includegraphics[width=1\linewidth]{"Figures/An Increase in the Investment Rate"}
%				\end{figure}
%			\end{column}%
%		\end{columns}
%\end{frame}

%\begin{frame}[wide]{Measuring the State of the Economy}
%	\begin{columns} 
%		% Column 1
%		\begin{column}{.5\textwidth}
%			\begin{table}[htbp]
%				\centering
%				\begin{tabular}{lrrrr}
%					& 2005  & 2006  & 2007  & 2008 \\
%					\hline
%					GDP   & 12.6  & 13.4  & 14.0    & 14.3 \\
%					(in trillions of \$) &       &       &       &  \\
%					Growth rate GDP &       & 0.063 & 0.045 & 0.021 \\
%					\hline
%				\end{tabular}%
%				\label{tab:addlabel}%
%			\end{table}%
%		\end{column}
%		% Column 2    
%		\begin{column}{.5\textwidth}			
%			
%			
%		\end{column}%
%	\end{columns}
%\end{frame}

\begin{frame}[wide]{Next Lecture}
	\begin{itemize}
		\item Next lecture will begin Ch.7
	\end{itemize}
\end{frame}

\end{document}